\hypertarget{python-3.12-native-generics-pep-695-type-statement-and-subinterpreters}{%
\chapter{Python 3.12: Native Generics (PEP 695), Type statement, and
Subinterpreters}\label{python-3.12-native-generics-pep-695-type-statement-and-subinterpreters}}

\hypertarget{pep-695-type-parameter-syntax}{%
\subsection{18.1 PEP 695 Type Parameter
Syntax}\label{pep-695-type-parameter-syntax}}

Python 3.12 introduced a clean, native syntax for generic classes,
generic functions, and type aliases using type parameters enclosed in
square brackets.

\begin{lstlisting}[language=Python]
# Generic Function
def reverse[T](items: list[T]) -> list[T]:
    return items[::-1]

# Generic Class
class Stack[T]:
    def __init__(self) -> None:
        self._data: list[T] = []

# Modern Type Alias
type Vector3D[T] = tuple[T, T, T]
\end{lstlisting}

\begin{center}\rule{0.5\linewidth}{0.5pt}\end{center}

\hypertarget{hidden-scopes-and-lazy-evaluation}{%
\subsection{18.2 Hidden Scopes and Lazy
Evaluation}\label{hidden-scopes-and-lazy-evaluation}}

Historically, generic class constraints (e.g.,
\passthrough{\lstinline!T = TypeVar('T', bound=int)!}) were evaluated
eagerly at module import time, which often led to circular import issues
and name resolution errors. Under PEP 695, CPython resolves this by
evaluating type parameter bounds and type aliases using \textbf{hidden
compiler-generated scopes}.

\hypertarget{scope-encapsulation}{%
\subsubsection{1. Scope Encapsulation}\label{scope-encapsulation}}

When the compiler encounters type parameters or the
\passthrough{\lstinline!type!} statement, it wraps the evaluation code
inside a nested lexical scope (similar to a hidden function block).

\begin{lstlisting}[language=Python]
# Conceptual translation of type alias Vector3D[T]
class Vector3D:
    # Under the hood, T is evaluated lazily inside a nested compiler scope
    def __lazy_type_eval__(T):
        return tuple[T, T, T]
\end{lstlisting}

\hypertarget{typealiastype-struct-representation}{%
\subsubsection{\texorpdfstring{2. \texttt{TypeAliasType} Struct
representation}{2. TypeAliasType Struct representation}}\label{typealiastype-struct-representation}}

At runtime, the \passthrough{\lstinline!type!} statement creates an
instance of \passthrough{\lstinline!typing.TypeAliasType!}. The CPython
class definition wraps: * \passthrough{\lstinline!\_\_name\_\_!}: Name
of the alias. * \passthrough{\lstinline!\_\_value\_\_!}: The aliased
type (evaluated lazily upon first access via its descriptor getter). *
\passthrough{\lstinline!\_\_type\_params\_\_!}: Tuple of type
parameters.

\begin{center}\rule{0.5\linewidth}{0.5pt}\end{center}
